
We design a matched-budget evaluation framework that enables direct method comparison across benchmarks. The key design principle: hold computational budget constant while varying method and benchmark.

\subsubsection{Overview}

We evaluate two sampling-based uncertainty methods---semantic entropy and self-consistency---under identical conditions:

\begin{table}[htbp]
\centering
\resizebox{\columnwidth}{!}{%
\begin{tabular}{lll}
\toprule
Parameter & Value & Rationale \\
\midrule
Sample count (N) & 10 & Standard in both original papers \\
Temperature & 0.7 & Enables diverse sampling \\
Max tokens & 256 & Sufficient for QA responses \\
Model & Llama-3-8B-Instruct & Open-weight, reproducible \\
Seed & 42 & Single seed for pilot \\
NLI Model & DeBERTa-large-mnli & For semantic clustering \\
Similarity & BERTScore (roberta-large) & For self-consistency \\
\bottomrule
\end{tabular}
}
\end{table}

\subsubsection{Semantic Entropy}

Semantic entropy quantifies uncertainty by clustering generated responses according to semantic equivalence, then computing entropy over the cluster distribution.

\textbf{Step 1: Generation.} For each query q, generate N=10 responses {r\_1, ..., r\_N} using temperature sampling.

\textbf{Step 2: Semantic Clustering.} Use DeBERTa-large-mnli to compute bidirectional NLI scores between response pairs. Responses r\_i and r\_j are considered semantically equivalent if both directions yield entailment probability above threshold $\tau =0.5$.

\textbf{Step 3: Entropy Computation.} Let {$C_1$, ..., $C_K$} be the resulting clusters. The cluster probability distribution is p\_k = |$C_k$|/N. Semantic entropy is H\_SE = -$\Sigma$ p\_k log p\_k.

Higher entropy indicates responses fall into many distinct semantic clusters, suggesting uncertainty.

\subsubsection{Self-Consistency}

Self-consistency measures agreement across generated responses using surface-level similarity.

\textbf{Step 1: Generation.} Identical to semantic entropy---N=10 responses per query.

\textbf{Step 2: Pairwise Similarity.} Compute BERTScore F1 between all response pairs: S\_ij = BERTScore-F1(r\_i, r\_j).

\textbf{Step 3: Consistency Score.} Self-consistency is the mean off-diagonal similarity: SC = (1/N(N-1)) $\Sigma _{i\neq j}$ S\_ij.

Higher consistency indicates similar responses. For hallucination detection, we use 1 - SC as the uncertainty score.

\subsubsection{Evaluation Protocol}

\textbf{Benchmarks.} TruthfulQA (817 questions total) and HaluEval-QA (10,000 samples total).

\textbf{Pilot Mode.} For pipeline validation, N=20 samples per dataset were used.

\textbf{Metrics.} Primary metric is AUROC with 95\% confidence intervals via bootstrap resampling (1,000 iterations).

\textbf{Gate Criterion.} AUROC > 0.55 as threshold for determining whether a method shows above-chance performance on a benchmark.
